%=====================================================================
\chapter{Acquisition, Wrangling and Cleaning}\label{ch:wrangling}
%=====================================================================
\locator{2}{Part II --- From Data to Insight}

\begin{outcomes}
By the end of this chapter you will be able to:
\begin{tight}
  \item \textbf{Select} an appropriate acquisition method and describe the risks each carries.
  \item \textbf{Perform} the standard cleaning operations: deduplication, type repair, missing-value handling and outlier treatment.
  \item \textbf{Choose} between join types and predict how each changes the row count.
  \item \textbf{Encode} categorical variables and \textbf{scale} numeric ones without corrupting the measurement scale.
  \item \textbf{Detect} and \textbf{prevent} data leakage --- the most expensive routine error in applied machine learning.
\end{tight}
\end{outcomes}

\begin{prereq}
\chref{data} (types, scales, quality dimensions, missingness mechanisms) and
\chref{lifecycle} (framing question 4, on what is knowable at prediction time).
\end{prereq}

\begin{keyterms}
ETL / ELT $\bullet$ API $\bullet$ web scraping $\bullet$ deduplication $\bullet$
imputation $\bullet$ join (inner, left, outer) $\bullet$ one-hot encoding $\bullet$
ordinal encoding $\bullet$ normalisation $\bullet$ standardisation $\bullet$
binning $\bullet$ data leakage $\bullet$ pipeline
\end{keyterms}

\section{Where Data Comes From}

\rowcolors{2}{white}{lightbg}
\begin{longtable}{@{}L{2.9cm}L{5.0cm}L{6.7cm}@{}}
\toprule
\rowcolor{primary!15}
\textbf{Source} & \textbf{Typical use} & \textbf{The risk you must manage} \\
\midrule
\endhead
Database query & Internal transactional systems & The production database is not a playground; a careless join can lock tables \\
File export (CSV, Excel) & One-off extracts, shared datasets & Silent type coercion --- Excel turns identifiers into dates and drops leading zeros \\
REST API & SaaS platforms, VLE data, threat feeds & Rate limits, pagination, and schemas that change without notice \\
Streaming (MQTT, Kafka) & IoT telemetry, live logs & Out-of-order and duplicate delivery; see \chref{dataeng} \\
Web scraping & Public pages with no API & Legal and ethical limits, brittle selectors, and terms of service \\
Manual entry / survey & Ground truth, labels, context & Human error, missingness that is rarely MCAR \\
\bottomrule
\end{longtable}

\begin{alertbox}[Before You Acquire Anything]
Ask two questions that are not technical: \emph{am I permitted to use this data
for this purpose}, and \emph{does it contain personal data}. Answering them after
building the model is how projects get cancelled at the last stage.
\chref{ethics} covers the substance; the point here is that acquisition is the
moment to ask.
\end{alertbox}

\section{The Cleaning Sequence}

Order matters. Doing these in the wrong sequence causes quiet corruption --- for
example, imputing before deduplicating lets duplicated rows vote twice on the
mean you impute with.

\dsfig{%
\begin{tikzpicture}[
  st/.style={rounded corners=3pt, draw=#1, line width=1pt, fill=#1!8, text=#1!50!black,
             text width=2.45cm, align=center, font=\scriptsize\bfseries, minimum height=1.0cm},
  arr/.style={-{Stealth[length=2.2mm]}, primary!55, line width=1pt},
  note/.style={font=\tiny\itshape, text=gray!55!black, align=center, text width=2.6cm}]

  \node[st=primary]   (a) at (0,0)     {1. Audit\\(\chref{data})};
  \node[st=secondary] (b) at (3.05,0)  {2. Fix types};
  \node[st=tealhl]    (c) at (6.10,0)  {3. Deduplicate};
  \node[st=greenhl]   (d) at (9.15,0)  {4. Join sources};
  \node[st=goldhl]    (e) at (12.2,0)  {5. Handle missing};
  \node[st=purplehl]  (f) at (3.05,-2.6){6. Treat outliers};
  \node[st=accent]    (g) at (6.10,-2.6){7. Encode \& scale};
  \node[st=primary]   (h) at (9.15,-2.6){8. Split\\train/test};

  \foreach \x/\y in {a/b,b/c,c/d,d/e}{\draw[arr] (\x) -- (\y);}
  \draw[arr] (e.south) .. controls +(0,-1.0) and +(0,1.0) .. (f.north);
  \draw[arr] (f) -- (g);
  \draw[arr] (g) -- (h);

  \node[note] at (3.05,-1.15) {deduplicate \emph{before} imputing, or duplicates skew the fill value};
  \node[note] at (12.2,-2.6) {\textbf{Steps 5--7 must be fitted on the training split only.}
                              Fitting them on everything is leakage.};
  \draw[-{Stealth[length=2mm]}, accent!70, dashed] (10.6,-2.6) -- (10.35,-2.6);
\end{tikzpicture}}%
{The cleaning sequence. The order is not arbitrary: each step assumes the previous one is
done. The note on the right is the single most important rule in the chapter and is
expanded in \S\ref{sec:leakage}.}{cleaning-sequence}

\subsection{Handling Missing Values}

\rowcolors{2}{white}{defbg}
\begin{longtable}{@{}L{3.3cm}L{4.6cm}L{6.6cm}@{}}
\toprule
\rowcolor{primary!15}
\textbf{Strategy} & \textbf{How} & \textbf{When it is appropriate} \\
\midrule
\endhead
Drop rows & Remove any row with a missing value & Missingness is MCAR and rare ($<5\%$). Otherwise you are deleting a biased subset \\
Drop column & Remove the feature entirely & More than roughly half missing and no reason to think absence is informative \\
Mean / median fill & Replace with the column's centre & Numeric, MCAR or MAR, and the variable is not central to the model \\
Forward fill & Carry the last observation forward & Time series with slow-changing values --- sensor telemetry especially \\
Model-based & Predict the missing value from other columns & Valuable feature, MAR, enough data to justify the complexity \\
Missing indicator & Add a boolean \texttt{was\_missing} column & \textbf{MNAR} --- the absence itself carries signal. Usually combined with a simple fill \\
\bottomrule
\end{longtable}

\begin{conceptbox}[The Default You Should Adopt]
When unsure, add the missing indicator \emph{and} fill with the median. It costs
one column, never destroys information, and lets the model decide whether absence
matters. In learning analytics and security logging this is almost always the
right call, because absence is rarely accidental (\chref{data}).
\end{conceptbox}

\section{Joining Data}

Real projects always need more than one table. The join type decides how many rows
survive, and choosing wrongly is a common source of silently wrong analyses.

\dsfig{%
\begin{tikzpicture}[font=\scriptsize]
\foreach \i/\name/\desc [count=\c from 0] in {%
  0/INNER/{only matching\\rows survive},
  1/LEFT/{all of A; nulls\\where B missing},
  2/OUTER/{everything from\\both sides}}{
  \pgfmathsetmacro{\xo}{\c*4.6}
  \begin{scope}[shift={(\xo,0)}]
    \ifnum\i=0
      \fill[gray!12] (-0.65,0) circle (1.0cm);
      \fill[gray!12] (0.65,0) circle (1.0cm);
      \begin{scope}\clip (-0.65,0) circle (1.0cm); \fill[primary!45] (0.65,0) circle (1.0cm);\end{scope}
    \fi
    \ifnum\i=1
      \fill[primary!45] (-0.65,0) circle (1.0cm);
      \fill[gray!12] (0.65,0) circle (1.0cm);
      \begin{scope}\clip (-0.65,0) circle (1.0cm); \fill[primary!45] (0.65,0) circle (1.0cm);\end{scope}
    \fi
    \ifnum\i=2
      \fill[primary!45] (-0.65,0) circle (1.0cm);
      \fill[primary!45] (0.65,0) circle (1.0cm);
    \fi
    \draw[primary, line width=1pt] (-0.65,0) circle (1.0cm);
    \draw[primary, line width=1pt] (0.65,0) circle (1.0cm);
    \node[font=\scriptsize\bfseries, text=primary!60!black] at (-1.35,1.15) {A};
    \node[font=\scriptsize\bfseries, text=primary!60!black] at (1.35,1.15) {B};
    \node[font=\scriptsize\bfseries, text=primary] at (0,-1.5) {\name};
    \node[align=center, text=gray!55!black, font=\tiny] at (0,-2.15) {\desc};
  \end{scope}}

\node[anchor=north, align=left, text width=13.5cm, font=\scriptsize, text=accent!70!black] at (4.6,-3.2)
  {\textbf{Always check the row count before and after a join.} If it \emph{grew},
   the key was not unique on one side and rows have been duplicated --- a
   many-to-many join is almost never what you intended.};
\end{tikzpicture}}%
{The three joins you will use. The default in most tools is inner, which silently discards
non-matching rows; if 300 of your 2{,}000 students have no VLE record, an inner join
removes them and your sample is no longer the cohort.}{join-types}

\section{Encoding and Scaling}

\begin{definitionbox}[One-Hot vs.\ Ordinal Encoding]
\term{One-hot encoding} turns a nominal variable with $k$ categories into $k$
binary columns, preserving the fact that no order exists.
\term{Ordinal encoding} maps categories to integers, and is correct \emph{only}
when the categories genuinely have an order (low $<$ medium $<$ high).
\end{definitionbox}

\begin{definitionbox}[Normalisation vs.\ Standardisation]
\term{Min--max normalisation} rescales to $[0,1]$:
$x' = \dfrac{x - \min}{\max - \min}$.\\[3pt]
\term{Standardisation} (z-score) centres and rescales by spread:
$z = \dfrac{x - \bar{x}}{s}$, giving mean 0 and standard deviation 1.
\end{definitionbox}

\begin{worked}[Scaling the Student Vectors]
From \chref{math}: $a = [92, 8, 12]$, $b = [45, 3, 0]$, with plausible ranges
attendance $[0,100]$, quizzes $[0,10]$, posts $[0,25]$.

\smallskip
\textbf{Unscaled distance:} $\sqrt{47^2 + 5^2 + 12^2} = \sqrt{2378} \approx 48.8$,
of which attendance contributed 93\%.

\textbf{Min--max normalised:}\\
$a' = [0.92,\; 0.80,\; 0.48]$, \quad $b' = [0.45,\; 0.30,\; 0.00]$.

\textbf{Scaled distance:}
$\sqrt{0.47^2 + 0.50^2 + 0.48^2} = \sqrt{0.2209 + 0.25 + 0.2304}
= \sqrt{0.7013} \approx 0.837$.

\smallskip
\textbf{Contribution now:} attendance 31\%, quizzes 36\%, posts 33\% --- roughly
equal, which is what ``these three features matter comparably'' actually means.
Note this is a \emph{modelling decision}, not a correction: you have declared the
three features equally important. If attendance really is more important, say so
with a weight, deliberately.
\end{worked}

\rowcolors{2}{white}{lightbg}
\begin{longtable}{@{}L{4.4cm}L{4.6cm}L{5.5cm}@{}}
\toprule
\rowcolor{primary!15}
\textbf{Method needs scaling?} & \textbf{Because} & \textbf{Examples} \\
\midrule
\endhead
\textbf{Yes --- essential} & Uses distances or gradients & k-NN, k-means, SVM, PCA, neural networks, ridge/lasso \\
\textbf{No --- indifferent} & Splits on one feature at a time, so units are irrelevant & Decision trees, random forests, gradient boosting \\
\bottomrule
\end{longtable}

\section{Data Leakage}\label{sec:leakage}

\begin{definitionbox}[Data Leakage]
\term{Leakage} occurs when information that would not be available at prediction
time influences training. The model scores brilliantly in testing and fails in
production, because it learned to use a fact from the future.
\end{definitionbox}

\dsfig{%
\begin{tikzpicture}[font=\scriptsize,
  blk/.style={rounded corners=3pt, minimum height=0.85cm, align=center, font=\scriptsize\bfseries},
  lab/.style={font=\tiny, text=gray!55!black, align=center}]

% ---- WRONG ----
\node[font=\scriptsize\bfseries, text=accent, anchor=west] at (0,2.35) {\faIcon{times-circle}~WRONG: scale first, then split};
\node[blk, fill=gray!18, text=gray!45!black, minimum width=7.2cm] (all1) at (3.6,1.6) {whole dataset};
\node[blk, fill=accent!22, text=accent!60!black, minimum width=7.2cm] (sc1) at (3.6,0.7)
     {compute mean \& sd from \emph{all} rows, scale};
\node[blk, fill=primary!18, text=primary!60!black, minimum width=5.0cm] (tr1) at (2.5,-0.25) {train};
\node[blk, fill=goldhl!25, text=goldhl!45!black, minimum width=2.0cm] (te1) at (6.2,-0.25) {test};
\draw[-{Stealth[length=1.8mm]}, gray!60] (all1) -- (sc1);
\draw[-{Stealth[length=1.8mm]}, gray!60] (sc1.south) -- (4.0,-0.25+0.43);
\node[lab, text=accent, text width=4.2cm, anchor=west] at (7.5,0.7)
     {the test set's mean has leaked into the training data};

% ---- RIGHT ----
\node[font=\scriptsize\bfseries, text=greenhl!55!black, anchor=west] at (0,-1.55)
     {\faIcon{check-circle}~RIGHT: split first, then fit the scaler on train only};
\node[blk, fill=gray!18, text=gray!45!black, minimum width=7.2cm] (all2) at (3.6,-2.3) {whole dataset};
\node[blk, fill=primary!18, text=primary!60!black, minimum width=5.0cm] (tr2) at (2.5,-3.25) {train};
\node[blk, fill=goldhl!25, text=goldhl!45!black, minimum width=2.0cm] (te2) at (6.2,-3.25) {test};
\node[blk, fill=greenhl!22, text=greenhl!45!black, minimum width=5.0cm] (sc2) at (2.5,-4.2)
     {fit scaler here};
\node[blk, fill=greenhl!10, text=greenhl!45!black, minimum width=2.0cm] (ap2) at (6.2,-4.2)
     {apply it here};
\draw[-{Stealth[length=1.8mm]}, gray!60] (all2.south) -- (tr2.north);
\draw[-{Stealth[length=1.8mm]}, gray!60] (all2.south) -- (te2.north);
\draw[-{Stealth[length=1.8mm]}, gray!60] (tr2) -- (sc2);
\draw[-{Stealth[length=1.8mm]}, greenhl!70] (sc2) -- (ap2);
\node[lab, text=greenhl!55!black, text width=4.2cm, anchor=west] at (7.5,-4.2)
     {the test set is never used to learn anything --- it only gets scored};
\end{tikzpicture}}%
{The most common leakage in student projects. Any transformation that \emph{learns}
something from the data --- scaling, imputation, encoding, feature selection --- must be
fitted on the training split alone and then applied to the test split.}{leakage-split}

\begin{pitfall}[Five Leaks and How to Spot Them]
\begin{tight}
  \item \textbf{Fitting a scaler or imputer on the full dataset.} Fix: fit on train,
        transform test. Use a \texttt{Pipeline} so it is impossible to forget.
  \item \textbf{A feature computed after the outcome.}
        \texttt{final\_attendance} to predict failure; \texttt{ticket\_close\_time}
        to predict overrun. Fix: framing question 4.
  \item \textbf{Random splitting of time-series data.} Training on Thursday to
        predict Wednesday. Fix: split chronologically (\chref{timeseries}).
  \item \textbf{Duplicate rows spanning the split.} The same student in both train
        and test through a merge error. Fix: split by \emph{entity}, not by row.
  \item \textbf{Feature selection on the whole dataset.} Choosing the top 20
        correlated features before splitting has already consulted the test set.
\end{tight}
\textbf{The universal warning sign:} a model that is far more accurate than the
problem plausibly allows. If your dropout classifier reaches 99\%, look for a
leak before celebrating.
\end{pitfall}

%---------------------------------------------------------------------
\begin{fourlenses}{Cleaning and Leakage}
\lensLA{Joining registry to VLE data is a \emph{left} join from registry, never an
inner join --- students with no VLE activity are precisely the at-risk group, and an
inner join deletes them. Classic leak: using \texttt{total\_submissions\_in\_semester}
to predict a week-4 outcome. Missing engagement data should be encoded with an
indicator, not imputed, because it is MNAR.}

\lensPM{Effort figures are self-reported and cluster on round numbers, so treat
values as ordinal-ish rather than precise. Classic leak: including
\texttt{actual\_completion\_date} or final story-point count when predicting
mid-sprint completion. Deduplicate carefully --- the same work item often appears
under two IDs after a project is re-scoped.}

\lensIOT{Forward fill is the natural imputation for telemetry, because a
temperature that was $70^\circ$ a second ago is probably still about $70^\circ$
--- but cap how far it carries, or a dead sensor will appear healthy for hours.
Classic leak: normalising the whole telemetry history including the future.
Deduplicate on (device ID, timestamp): at-least-once delivery guarantees duplicates.}

\lensSEC{Do \emph{not} remove outliers --- they are the attacks. Encode
categorical fields such as protocol and port with one-hot, never ordinal, or the
model will infer that port 443 is ``greater than'' port 80. Classic leak: features
derived from the incident report, which by definition exists only after the attack
was detected. Split chronologically, since a model trained on future attack
patterns to detect past ones proves nothing.}
\end{fourlenses}

%---------------------------------------------------------------------
\begin{lab}[A Leak-Proof Cleaning Pipeline]
The single best habit in applied machine learning: put every learned
transformation inside a \texttt{Pipeline}, so it cannot be fitted on the test set
by accident.

\begin{lstlisting}[language=Python]
import pandas as pd
from sklearn.model_selection import train_test_split
from sklearn.pipeline import Pipeline
from sklearn.compose import ColumnTransformer
from sklearn.impute import SimpleImputer
from sklearn.preprocessing import StandardScaler, OneHotEncoder

df = pd.read_csv("students.csv")
y  = df.pop("failed")

numeric     = ["attendance_pct", "quizzes_done", "forum_posts"]
categorical = ["programme", "entry_route"]

# 1. SPLIT FIRST -- before anything is learned from the data.
X_tr, X_te, y_tr, y_te = train_test_split(
    df, y, test_size=0.2, stratify=y, random_state=42)

# 2. Declare transformations; nothing is fitted yet.
prep = ColumnTransformer([
    ("num", Pipeline([("impute", SimpleImputer(strategy="median",
                                               add_indicator=True)),   # MNAR-safe
                      ("scale",  StandardScaler())]), numeric),
    ("cat", OneHotEncoder(handle_unknown="ignore"), categorical),
])

# 3. fit_transform on TRAIN, transform only on TEST.
X_tr_ready = prep.fit_transform(X_tr)
X_te_ready = prep.transform(X_te)      # <- .transform, never .fit_transform

print(X_tr.shape, "->", X_tr_ready.shape)   # wider: one-hot + indicators
\end{lstlisting}

\textbf{Check your understanding:} what goes wrong if line 3's last call is
changed to \texttt{fit\_transform}? Name the leak and the symptom you would see.
\end{lab}

%---------------------------------------------------------------------
\begin{checkpoint}
\begin{tightnum}
  \item You join 2{,}000 student rows to a VLE table and get 2{,}380 rows back.
        What has happened, and what should you check?
  \item Which of these need scaling before use: k-NN, random forest, k-means,
        gradient boosting?
  \item A model predicting 30-day equipment failure achieves 99.8\% accuracy on
        held-out data. State the first three things you would check.
\end{tightnum}
\end{checkpoint}

%---------------------------------------------------------------------
\begin{chaptersummary}
\begin{tight}
  \item Acquisition method determines the failure modes you inherit; check
        permission and personal-data status \emph{before} building.
  \item Clean in order: audit, types, deduplicate, join, missing, outliers, encode
        and scale, split.
  \item Under MNAR, add a missing indicator rather than imputing away the signal.
  \item Inner joins silently delete non-matching rows; always compare row counts
        before and after.
  \item One-hot for nominal, ordinal encoding only for genuinely ordered
        categories. Scale for distance- and gradient-based methods; trees do not
        care.
  \item Leakage is the defining error of the field. Split first, fit
        transformations on train only, and treat implausibly good results as a
        symptom rather than a success.
\end{tight}
\end{chaptersummary}

\begin{reviewq}
  \item \textbf{(MCQ)} Which method is indifferent to feature scaling?
        (a)~k-nearest neighbours (b)~k-means (c)~random forest (d)~support vector machine.
  \item \textbf{(MCQ)} Fitting a \texttt{StandardScaler} on the full dataset before
        splitting causes: (a)~underfitting (b)~leakage (c)~class imbalance
        (d)~multicollinearity.
  \item \textbf{(Short)} Explain when ordinal encoding is correct and give an example
        where using it would be a serious error.
  \item \textbf{(Short)} Why must forward fill be capped when imputing sensor
        telemetry?
  \item \textbf{(Short)} Describe two situations in which an inner join produces a
        biased dataset without producing any error message.
  \item \textbf{(Applied)} You must predict, at the end of week 4, which students
        will fail. From this list, mark each feature usable or leaking and justify:
        \texttt{weeks1\_4\_attendance}, \texttt{final\_exam\_mark},
        \texttt{n\_submissions\_by\_week4}, \texttt{total\_semester\_logins},
        \texttt{entry\_qualification}, \texttt{withdrew\_date}.
  \item \textbf{(Applied)} Write the cleaning sequence you would apply to an IoT
        dataset of 400 motors reporting every 10 seconds for a year, stating for
        each step the specific decision you would take and why.
\end{reviewq}
